% D.62; independent root and FGH full proof audits passed.
\subsection{A nonzero Jannsen kernel for every Alexander twist}

\begin{statement}{D.62}{Proved}{Infinite ramification never removes the nonseparated second cohomology}
Let \(s=(a,b)\) be any completed abelian Kummer pair over
\(\mathbb Q\), and let \(A_s=\Lambda C\) have the actual
Alexander action of D.35. No finite ramification hypothesis
and no metabelian lifting hypothesis are imposed. Set
\(W_r=A_s/\mathfrak m^r A_s\), where
\(\mathfrak m=(5,u,v)\). Then
\[
 \varprojlim\nolimits_r^1 H^1(G_{\mathbb Q},W_r)\ne0.
                                                        \tag{D.62.1}
\]
Consequently the continuous-cohomology comparison fits into
the exact sequence with nonzero first term
\[
 0\longrightarrow
 \varprojlim\nolimits_r^1 H^1(G_{\mathbb Q},W_r)
 \longrightarrow H^2_{\mathrm{cts}}(G_{\mathbb Q},A_s)
 \longrightarrow\varprojlim_r H^2(G_{\mathbb Q},W_r)
 \longrightarrow0.
                                                        \tag{D.62.2}
\]
With the continuous-cochain quotient topology defined in
the proof, the first term is exactly the closure of zero.
In particular this \(H^2\) is not Hausdorff for any twist.
Every class in that first term restricts to zero in
\(H^2_{\mathrm{cts}}(G_{\mathbb Q_v},A_s)\) at every place
\(v\).

This proves nonvanishing of the unrestricted derived-limit
term even for infinitely ramified pairs, including the
rank-zero modules of D.60--D.61. It does not evaluate its
isomorphism type or the entire second cohomology for every
such pair.
\end{statement}

\begin{proof}
Write \(B_r=H^1(G_{\mathbb Q},W_r)\). Each \(B_r\)
is countable. Indeed there are countably many finite
Galois extensions of \(\mathbb Q\), and every continuous
cocycle with finite coefficients, together with its finite
action, factors through one of these finite quotients.

\paragraph{Every finite stage has infinitely many surviving
residues at the first coefficient level.}
Fix \(r\). Choose a finite set \(S\) containing five and
infinity such that the finite module \(W_r\) is unramified
outside \(S\). Such an \(S\) always exists, regardless of
the ramification of the full compact module \(A_s\).
Let \(X=\operatorname{Spec}\mathbb Z[S^{-1}]\), with
infinity omitted from the ring notation. Finite-coefficient
localization gives the exact segment
\[
 H^1(G_{\mathbb Q},W_r)
 \xrightarrow{\partial_r}
 \bigoplus_{q\notin S}H^0(\mathbb F_q,W_r(-1))
 \xrightarrow{\delta_r} H^2(G_{\mathbb Q,S},W_r).
                                                        \tag{D.62.3}
\]
For completeness, apply localization on \(X\) after
removing a finite set of closed points. At an unramified
prime, purity identifies the cohomology with support
with the residue-field cohomology shifted by two and
twisted by minus one. Taking the filtered union over
the removed finite sets gives (D.62.3), since continuous
finite-valued cochains factor through finite Galois
quotients. The comparison of the etale and Galois groups
is \cite[II, Proposition 2.9]{DE_MilneADT}. This is the
finite-coefficient calculation of (D.48.4); no inverse
limit is taken with the current \(S\) fixed across \(r\).
The last group in (D.62.3) is finite: its kernel of
localization is finite by
\cite[I, Theorem 4.10(a)]{DE_MilneADT}, and its image
lies in the finite product of the finite local second
cohomology groups at \(v\in S\)
\cite[I, Corollary 2.3]{DE_MilneADT}.

The exponent of \(W_r\) divides \(5^r\). Take a finite
Galois extension cutting out both its actual action
\(\rho_r\) and \(\mu_{5^r}\). Infinitely many primes
\(q\notin S\) split completely in this extension, by
Chebotarev \cite[Chapter V, Theorem 3.23]{PBC_MilneCFT}. At each such prime,
\[
 \rho_r(\operatorname{Frob}_q)=1,
 \qquad q\equiv1\pmod{5^r},
 \qquad H^0(\mathbb F_q,W_r(-1))=W_r(-1).
                                                        \tag{D.62.4}
\]
Choose at each such residue field a compatible basis
for the Tate twist and let \(w_q\) be the residue
corresponding to \(C\bmod\mathfrak m^r\). Its reduction
in \(W_1(-1)\) is nonzero.

Because the target of \(\delta_r\) is finite, an infinite
subset of these primes has one common value
\(\delta_r(w_q)\). Fix a prime \(q_*\) in that subset.
For every other \(q\) there, the element with residue
\(w_q\) at \(q\), residue \(-w_{q_*}\) at \(q_*\),
and zero at all other primes lies in \(\ker\delta_r\).
By (D.62.3) it is the residue of a class in \(B_r\).
Reduction to \(B_1\) leaves a nonzero residue at its
own prime \(q\) and none at the other chosen primes
except the common \(q_*\). These reduced classes are
linearly independent over \(\mathbb F_5\), as their
individual residues show. Consequently
\[
 \dim_{\mathbb F_5}\operatorname{im}(B_r\longrightarrow B_1)
                              =\aleph_0
                   \quad\hbox{for every }r\ge1.
                                                        \tag{D.62.5}
\]
The upper bound follows from the countability already
proved. This uses the actual finite action, not only
the cyclotomic character.

\paragraph{Integral cocycles have a uniformly finite first image.}
The groups \(H^0(G_{\mathbb Q},W_r)\) are finite, so
their inverse system is Mittag--Leffler. The continuous
cochain comparison therefore gives
\[
 H^1_{\mathrm{cts}}(G_{\mathbb Q},A_s)
                         \simeq\varprojlim_r B_r.
                                                        \tag{D.62.6}
\]
One may derive this directly from the cochain complexes:
\(C^\bullet(G_{\mathbb Q},A_s)=
\varprojlim_r C^\bullet(G_{\mathbb Q},W_r)\), and each
degreewise transition is surjective, since a continuous
finite-valued cochain can be lifted value by value.
This also proves (D.62.2), or one may apply
\cite[Proposition 1.6 and Theorem 2.2]{DE_Jannsen1988}.

For \(r=1\), the set \(S_r\) of (D.47.1) is just
\(\{5,\infty\}\): every Kummer valuation is integral,
so none has negative five-adic valuation. Thus every
integral \(A_s\)-cocycle, reduced to \(W_1\), factors
through \(G_{\mathbb Q,\{5,\infty\}}\). The module
\(W_1\simeq\mathbb F_5(2)\) is independent of the
abelian twist. Its finite-ramification first cohomology
is finite. In particular
\[
 \operatorname{im}\bigl(\varprojlim_r B_r\longrightarrow B_1\bigr)
                             \text{ is finite}.
                                                        \tag{D.62.7}
\]
The finiteness, rather than any vanishing calculation
for this group, is all that is needed.

If the countable tower \((B_r)\) were Mittag--Leffler,
let \(B'_r\) be the eventual image of \(B_j\) in
\(B_r\). Its transition maps would be surjective:
choose \(j\) large enough that both eventual images
have stabilized and lift from \(B_j\).
Recursive lifting then makes
\(\varprojlim_r B'_r\to B'_1\) surjective, and the
limits of \((B'_r)\) and \((B_r)\) agree. But by
(D.62.5), \(B'_1\) would be infinite-dimensional,
contradicting (D.62.7). Hence \((B_r)\) is not
Mittag--Leffler. For countable abelian groups the
converse to the usual derived-limit vanishing criterion
holds: \(\varprojlim^1 B_r=0\) would imply the
Mittag--Leffler condition. This is the Gray criterion
stated in \cite[Lemma 4.15, printed p.229]{DE_Jannsen1988}.
It proves (D.62.1).

\paragraph{The topology and the local invisibility.}
Give continuous \(A_s\)-valued cochains the uniform
\(\mathfrak m\)-adic topology: a basic neighborhood
of zero consists of cochains taking all their values
in \(\mathfrak m^rA_s\). Give the two-cocycle group
the subspace topology and \(H^2\) the quotient by
continuous coboundaries. Suppose a two-cocycle \(z\)
has zero class at every finite coefficient level.
For each \(r\), choose a continuous finite-valued
one-cochain \(b_r\) with
\(z\bmod\mathfrak m^r=db_r\), and lift its finitely
many values to \(A_s\). The resulting continuous
one-cochain \(\widetilde b_r\) satisfies
\[
                  z-d\widetilde b_r
                       \in C^2(G_{\mathbb Q},\mathfrak m^rA_s).
                                                        \tag{D.62.8}
\]
Thus \(z\) is in the closure of the coboundaries.
Conversely an approximating coboundary to this precision
makes the reduced class zero. The kernel in (D.62.2)
is therefore exactly the closure of zero for the
displayed quotient topology. Its nonzero size proves
non-Hausdorffness.

At every finite place \(v\), the groups
\(H^1(G_{\mathbb Q_v},W_r)\) are finite
\cite[I, Corollary 2.3]{DE_MilneADT}; at the real place
the analogous statement is immediate. Their derived
limit vanishes. The local analogue of (D.62.2) is
therefore injective into its inverse limit of finite
coefficient second cohomology. Restriction of a global
class in the first term of (D.62.2) is zero at every
such finite level, and consequently is zero locally.
This proves the final assertion without claiming a
global Hasse principle for the compact module.
\end{proof}
